\section{Notation and Exact State}
\label{app:state}

\subsection{Eligible exposure and frontiers}

At a fixed routing state $t$, reuse the credit-eligible set
$\Omega_t(s)=\{x\mid q_t(s,x)\text{ exists and }x\notin\mathcal B_s\}$ from
\Cref{sec:setting}.  The frontier is
\[
\F_t(x)=\argmax_{\substack{s\\x\in\Omega_t(s)}}q_t(s,x),
\]
including all ties, even when all observed scores are zero.  For each instance, define
\[
\mathcal C_t(x)=\{u\mid x\in\Omega_t(u)\},\quad N_x=|\mathcal C_t(x)|,\quad
y_{sx}=\ind\{s\in\F_t(x)\},\quad K_x=\sum_{u\in\mathcal C_t(x)}y_{ux}.
\]
When $N_x\ge2$, the per-instance credit is
\[
r_{sx}=y_{sx}-\frac{K_x-y_{sx}}{N_x-1}.
\]
The relative credit $G_s$ averages these per-instance credits.
\[
\begin{gathered}
\Omega_t^{\mathrm{cmp}}(s)=\{x\in\Omega_t(s)\mid N_x\ge2\},\qquad
E_s^{\mathrm{cmp}}=|\Omega_t^{\mathrm{cmp}}(s)|,\\
R_s=\sum_{x\in\Omega_t^{\mathrm{cmp}}(s)}r_{sx},\qquad
G_s=R_s/E_s^{\mathrm{cmp}}.
\end{gathered}
\]
An $N_x=1$ observation remains committed but does not enter the aggregate.
Whenever a new finite observation changes $\mathcal C_t(x)$ or $\F_t(x)$, the
router recomputes $N_x$, $K_x$, and every affected candidate contribution on
that instance.  Candidates with $E_s^{\mathrm{cmp}}=0$ have undefined credit.

\subsection{Inherited proposal history}

If parent $s$ is proposed on batch $P$, the prospective child cannot use its
current proposal instances or any ancestor birth-proposal instances as
admission evidence.  Let $P_a^{\mathrm{birth}}$ be the recorded proposal batch
that created accepted candidate $a$, with $P_0^{\mathrm{birth}}=\varnothing$
for the seed.  The child's proposal history is
\[
X(s,P)=P\cup\bigcup_{a\in\operatorname{Anc}(s)\cup\{s\}}
P_a^{\mathrm{birth}}.
\]
Proposal history records the instances used in a candidate's own and its
ancestors' proposal batches.  Admission observations retain their role as evaluation evidence.
The proposal history $\mathcal B_s$
governs which committed observations may contribute to candidate $s$'s relative
credit, and $X(s,P)$ extends that exclusion to the prospective child's admission.

\section{Algorithmic Details}
\label{app:algorithm}

\subsection{Proposal allocation implementation}

\Cref{eq:repairable_gap,eq:proposal_assignment} define the
proposal objective.  The implementation builds the complete slot--batch
weight matrix from the committed cache and solves the resulting
maximum-weight assignment exactly.  Duplicate parent labels remain distinct
left-side slot vertices.  Missing cache entries remain unchanged during assignment.

\subsection{Direct legal admission implementation}

For each of the $m\le w$ remaining proposal tasks, the benchmark evaluator samples
one batch $A_i$ from that task's admission sampler after removing every
instance in $X(s_i,P_i)$.  For each instance in the sampled batch, the evaluator
selects the frontier candidate with defined, maximal
$G$ as reference, using seeded randomness for an exact tie and the sampled
parent when no frontier candidate has defined credit.  The resulting plan is hidden from the
generator, required missing reference observations are filled on the existing
reference programs, and the fixed child is admitted exactly when the paired
aggregate in \Cref{eq:admission_gate} is positive.

\subsection{Wave execution}

\begin{algorithm}[H]
\caption{One \method{} wave}
\label{alg:compass_wave}
\begin{algorithmic}[1]
\Require committed routing state $\mathcal H_t$, wave width $w$, active cutoff $\rho$
\Ensure updated committed state $\mathcal H_{t+1}$
\State Recompute every defined $G_s$ from the committed cache
\State $\A_t \gets$ tie-inclusive Top-$\rho$ candidates, sorted by $G_s$ and candidate index
\State If $\A_t=\varnothing$, set $\A_t\gets\{0\}$ \Comment{Seed fallback}
\State $(z_i)_{i\in\mathcal I_w} \gets \text{next $w$ parents from $\A_t$, continuing from the saved position}$
\State Draw $(P_j)_{j\in\mathcal I_w}$ and form $M_{ij}\gets H(z_i,P_j)$
\State $Y^\star \gets \Call{MaxWeightAssignment}{M}$
\ForAll{assigned parent--proposal pairs $(s_i,P_i)$ \textbf{in parallel}}
  \State Evaluate $s_i$ on $P_i$ and construct its reflection data
  \State Update the cache with new pairs or strictly higher parent scores
\EndFor
\State $\mathcal T\gets$ tasks with usable traces and reflection data, after configured perfect-score skipping
\State $m\gets|\mathcal T|\le w$
\ForAll{$(s_i,P_i)\in\mathcal T$ \textbf{in parallel}}
  \State $A_i\gets\Call{SampleAdmission}{s_i,P_i}$ subject to $A_i\cap X(s_i,P_i)=\varnothing$
  \State $r_i\gets\Call{BindReferences}{A_i}$, then freeze and hide $(A_i,r_i)$ from the generator
  \State Evaluate missing references and store their scores and outputs
  \State $p_i\gets\Call{Reflect}{s_i,P_i}$, then evaluate $p_i$ on $A_i$
  \State $\Delta_i\gets\sum_{x\in A_i}[q(p_i,x)-q(r_i(x),x)]$
  \If{$\Delta_i>0$}
    \State Commit $p_i$ and its lineage-eligible admission evidence
  \EndIf
\EndFor
\State Update the logical budget, frontiers, credits, and queue cursor, then persist $\mathcal H_{t+1}$
\end{algorithmic}
\end{algorithm}

Repeat \Cref{alg:compass_wave} while the cumulative logical call count is below
$B$, and complete the final wave before stopping.  Report both the target $B$
and actual calls.  Request-level retries add no logical task result.

\section{Properties of Admission, Credit, and Selection}
\label{app:proofs}

We prove the exclusion and paired-concentration properties of admission,
instance-wise credit conservation, finite-pool selection bounds, coverage of
the active set, and optimality of the proposal assignment.

\subsection{Why proposal and admission evidence are separate}

\begin{proof}[Proof of \Cref{prop:self_cert_leakage}]
Let the task distribution be uniform on a domain of size $K\ge n/\epsilon$.
After observing $S=(x_1,\ldots,x_n)$, let the proposal rule construct a program
whose reward is one exactly on the distinct elements appearing in $S$ and zero
elsewhere.  Every reused evaluation example belongs to that set, so the
empirical reward on $S$ is one.  The population reward is the fraction of the
domain covered by the observed set, at most $n/K\le\epsilon$.
\end{proof}

The construction motivates admission on examples outside those that shaped
the candidate.  Recording proposal history extends this separation through the
lineage, because a descendant can retain behavior adapted to an ancestor's
proposal batch.

\subsection{Paired admission}

\begin{theorem}[Conditional concentration for paired admission]
\label{thm:paired_admission}
Let $\mathcal V$ be the information visible to the generator when child $p$ is
fixed from parent $s$ and proposal batch $P$, and suppose that $p$ is
$\mathcal V$-measurable.  Before generation, the evaluator fixes an admission batch
$A=(X_1,\ldots,X_n)$, with $n\ge1$, outside $X(s,P)$ and binds each reference $r(X_j)$ from
the committed pre-child state and seeded tie-breaking randomness.  Admission
identities, references, and outcomes are not visible in $\mathcal V$.  Then
\[
A\cap P=\varnothing,
\qquad
A\cap P_a^{\mathrm{birth}}=\varnothing
\quad\text{for every }a\in\operatorname{Anc}(s)\cup\{s\}.
\]
After $p$ is fixed, reveal the paired outcomes sequentially.  Let
$(\mathcal H_j)_{j=0}^n$ be the filtration generated by the fixed child,
the frozen admission plan, and the first $j$ revealed paired outcomes.
Define
\[
Z_j=q(p,X_j)-q(r(X_j),X_j),
\qquad
\mu_j=\E[Z_j\mid\mathcal H_{j-1}],
\]
\[
\widehat\mu_A(p,r)=\frac1n\sum_{j=1}^n Z_j,
\qquad
\overline\mu_A(p,r)=\frac1n\sum_{j=1}^n\mu_j.
\]
Because $Z_j\in[-1,1]$, for every $\delta\in(0,1)$, with probability at
least $1-\delta$,
\begin{equation}
\left|\widehat\mu_A(p,r)-\overline\mu_A(p,r)\right|
\le
\sqrt{\frac{2\log(2/\delta)}{n}}.
\label{eq:paired_concentration}
\end{equation}
For the $i$th task, the implemented gate is equivalently
$\operatorname{admit}(p_i)\iff
\widehat\mu_{A_i}(p_i,r_i)>0$.
\end{theorem}

\begin{proof}[Proof of \Cref{thm:paired_admission}]
Because $X(s,P)$ contains both $P$ and every
$P_a^{\mathrm{birth}}$ for $a\in\operatorname{Anc}(s)\cup\{s\}$, legality
$A\cap X(s,P)=\varnothing$ gives the displayed disjointness relations.
After the child is fixed, each paired difference $Z_j$ lies in an interval of
length two.  Conditional Hoeffding--Azuma applied to
$Z_j-\E[Z_j\mid\mathcal H_{j-1}]$ gives
\[
\Pr\!\left(
\left|\widehat\mu_A(p,r)-\overline\mu_A(p,r)\right|\ge\epsilon
\right)
\le 2\exp\!\left(-\frac{n\epsilon^2}{2}\right).
\]
Setting the right-hand side to $\delta$ proves
\Cref{eq:paired_concentration}.
\end{proof}



\subsection{Per-instance credit sums to zero}

\begin{proof}[Proof of \Cref{prop:credit_conservation}]
For a fixed comparable instance $x$, use
$\sum_{s\in\mathcal C_t(x)}y_{sx}=K_x$ to obtain
\[
\begin{aligned}
\sum_{s\in\mathcal C_t(x)}r_{sx}
&=K_x-
\frac{\sum_{s\in\mathcal C_t(x)}(K_x-y_{sx})}{N_x-1}\\
&=K_x-\frac{N_xK_x-K_x}{N_x-1}=0.
\end{aligned}
\]
The statement is per instance, and the candidate averages $G_s$ need not sum to
zero because their comparable exposure sets and denominators can differ.
\end{proof}



\subsection{Finite-pool selection}

\begin{proposition}[Finite-pool mean selection]
\label{prop:bandwidth_precision}
For any $M,n\in\mathbb N_+$ and $\delta\in(0,1)$, let
$p_1,\ldots,p_M$ be adaptively generated candidates.  Let
$\mathcal J_{i-1}$ contain the history before candidate $p_i$ is evaluated,
with $p_i$ measurable with respect to that history.  Conditional on
$\mathcal J_{i-1}$, suppose $p_i$ receives $n$ independent and identically
distributed rewards in $[0,1]$ under a common benchmark evaluation
distribution, with mean $R(p_i)$.  Let $\widehat R_n(p_i)$ be their
arithmetic mean and define
\[
\widetilde p_{M,n}\in\argmax_{p\in\{p_1,\ldots,p_M\}}\widehat R_n(p).
\]
Then, with probability at least $1-\delta$,
\begin{equation}
R^\star-R(\widetilde p_{M,n})
\le
\underbrace{R^\star-\max_{i\le M}R(p_i)}_{\text{search gap}}
+\underbrace{2\sqrt{\frac{\log(2M/\delta)}{2n}}}_{\text{selection noise}}.
\label{eq:bandwidth_variance_regret}
\end{equation}
\end{proposition}

\begin{proof}[Proof of \Cref{prop:bandwidth_precision}]
Conditional on $\mathcal J_{i-1}$, Hoeffding's inequality gives
\[
\Pr\!\left(
|\widehat R_n(p_i)-R(p_i)|>\epsilon
\mid\mathcal J_{i-1}\right)\le2e^{-2n\epsilon^2}.
\]
Taking expectations over the history and a union bound over the $M$
candidates gives, with probability at least $1-\delta$,
\[
\max_{i\le M}|\widehat R_n(p_i)-R(p_i)|\le\epsilon,
\qquad \epsilon=\sqrt{\frac{\log(2M/\delta)}{2n}}.
\]
Let $p_M^\star$ maximize population utility in the pool.  On this event,
\[
R(p_M^\star)-R(\widetilde p_{M,n})
\le\widehat R_n(p_M^\star)-\widehat R_n(\widetilde p_{M,n})
+2\epsilon\le2\epsilon.
\]
Adding the search gap proves \Cref{eq:bandwidth_variance_regret}.
\end{proof}

Taking expectations and bounding regret on the failure event by one also
gives
\[
\E\!\left[R^\star-R(\widetilde p_{M,n})\right]
\le\E\!\left[R^\star-\max_{i\le M}R(p_i)\right]
+2\sqrt{\frac{\log(2M/\delta)}{2n}}+\delta.
\]
Along a nested candidate sequence, the search gap is nonincreasing in $M$, while
the selection-noise term decreases with $n$ for fixed $M$.

\paragraph{Selection by relative credit.}
\begin{corollary}[Credit-based selection]
\label{cor:g_selection}
Under the conditions of \Cref{prop:bandwidth_precision}, let
$\widehat p_G$ be a candidate with the largest defined $G$ in the pool and
let $\xi_G$ be as in \Cref{eq:g_selection_gap}.  Then
\Cref{eq:g_selection_regret} holds with probability at least $1-\delta$.
Moreover,
\[
\E\!\left[R^\star-R(\widehat p_G)\right]
\le\E\!\left[R^\star-\max_{i\le M}R(p_i)\right]
+2\sqrt{\frac{\log(2M/\delta)}{2n}}+\E[\xi_G]+\delta.
\]
\end{corollary}

\begin{proof}[Proof of \Cref{cor:g_selection}]
On the simultaneous concentration event used above,
\[
\begin{aligned}
R(p_M^\star)-R(\widehat p_G)
&\le\widehat R_n(p_M^\star)-\widehat R_n(\widehat p_G)+2\epsilon\\
&\le\xi_G+2\epsilon.
\end{aligned}
\]
Adding the search gap yields \Cref{eq:g_selection_regret}.  The expectation
bound follows because regret is at most one on the failure event and the
search gap and $\xi_G$ are nonnegative.
\end{proof}

\paragraph{Coinciding rankings for binary scores.}
Suppose the $M\ge2$ candidates share the same $n$ legal instances, each
candidate--instance pair has one binary score, and these are all their
comparable observations.  Thus $N_x=M$ and
$\widehat R_n(p_s)=n^{-1}\sum_x q(s,x)$.  On each instance, either at least
one score is one and $y_{sx}=q(s,x)$, or every score is zero and every
candidate is on the frontier.  In both cases,
\[
r_{sx}=q(s,x)-\frac{1}{M-1}\sum_{u\ne s}q(u,x).
\]
Averaging over the common instances gives
\begin{equation}
G_s=\frac{M}{M-1}\left(
\widehat R_n(p_s)-\frac1M\sum_{u=1}^M\widehat R_n(p_u)
\right).
\label{eq:binary_credit_mean}
\end{equation}
This is a shared positive affine transformation of the sample means, so
their rankings and ties coincide and $\xi_G=0$.

\subsection{Structural properties of the router}

\begin{proposition}[Proposal turns for a fixed active set]
For a fixed ordered active set of size $m$, every block of $m$ consecutive
queue draws contains each active candidate exactly once, including when a
width-$w$ wave wraps around the queue.
\end{proposition}

\begin{proof}
Index the fixed active order by $0,\ldots,m-1$ and let the persistent cursor
be $c$.  Consecutive draws are $(c+r)\bmod m$.  For
$r=0,\ldots,m-1$, these residues are pairwise distinct and cover every index.
Grouping the same sequence into blocks of width $w$ does not alter the argument,
including a block that wraps around the end of the queue.
\end{proof}

\begin{proposition}[Maximum-weight proposal matching]
Conditional on the $w$ drawn parent slots and $w$ proposal minibatches,
\Cref{eq:proposal_assignment} maximizes total matching weight over all
one-to-one slot--minibatch assignments.
\end{proposition}

\begin{proof}[Proof of maximum-weight proposal matching]
The $w$ parent slots and $w$ minibatches define a complete bipartite graph
with edge weight $H(z_i,P_j)$.  A feasible binary matrix in
\Cref{eq:proposal_assignment} is exactly a perfect matching, or equivalently a
permutation matrix.  Selecting the maximum-weight feasible matrix therefore
returns the maximum-weight matching.  Repeated parent labels do not change the graph because
their slot vertices remain distinct.
\end{proof}

\section{Detailed Experimental Protocol}
\label{app:protocol}

\subsection{Method configuration}
\label{app:frozen_identities}

The method in \Cref{sec:method} uses
relative credit, applies the tie-inclusive Top-$\rho$ cutoff
to all candidates with defined credit, uses the exact proposal assignment in
\Cref{eq:proposal_assignment}, and samples admission batches independently for
each remaining proposal task.  The implementation names this credit
\texttt{leave\_one\_out\_frontier\_credit}.  The protocol evaluates the selected
program once.

The Qwen3-8B runs select candidates 55, 78, 34, and 30 for HotPotQA,
IFBench, PUPA, and LiveBench-Math, respectively.  GEPA's
GPT-5.4-nano AIME-2025 result uses its first completed official evaluation.
Each saved run records rollout counts, candidate indices, tokens, and costs.
Its checkpoint retains lineage, queue, assignment, and cache updates for
analysis of the search mechanism.

\subsection{Official evaluation protocol}

Reported results use the candidate fixed at the end of optimization and a
complete official evaluation.  Each run records the selected program, source
version, configuration, and model.  IFBench requires 294 finite example scores
whose mean equals the aggregate.  AIME requires all 150 benchmark
problems.  HoVer requires 300 finite binary scores, and LiveBench-Math requires
all 126 frozen test positions under its official aggregation.

Evaluation uses the selected candidate with its matching implementation and
configuration.  Service failures are recorded separately from benchmark scores.
Preflight runs and intermediate checkpoints remain separate from the final
evaluation, and the candidate is fixed before its test scores are observed.

\subsection{Resource accounting}

Logical task calls and final-wave overshoot are reported separately from physical
language-model calls.  The latter include input/output tokens and route-priced
cost.  A local task deadline can produce a benchmark failure score when that
outcome is defined by the evaluator.  Other provider, schema, serialization,
and engineering errors terminate the evaluation
batch.  A resumable run must preserve the exact source, config, cache,
checkpoint, and random-number-generator state.
